\begin{figure}[H]
\centering
\resizebox{\linewidth}{!}{%
\begin{tikzpicture}[font=\scriptsize,
  filtro/.style={caja, text width=1.75cm, minimum height=1.5cm, align=center, inner sep=3pt}]

\node[filtro, draw=ambar, line width=1pt] (in) at (0,0) {\textbf{trama}\\[1pt]{\tiny MQTT · HTTP}};

\foreach \i/\nom/\det in {
  1/{autenticar}/{TLS + credencial},
  2/{analizar}/{¿se puede leer?},
  3/{validar}/{contra el registro},
  4/{normalizar}/{al modelo interno},
  5/{supervisar}/{¿quién calla?},
  6/{alertar}/{umbrales y caídas}}{
  \pgfmathsetmacro\x{\i*2.35}
  \node[filtro] (f\i) at (\x,0) {\textbf{\i}\\ \textbf{\nom}\\[1pt]{\tiny \det}};
}
\node[filtro, draw=verdeacento, line width=1.1pt, fill=verdeacento!6] (db) at (16.45,0)
  {\textbf{TimescaleDB}\\[1pt]{\tiny series de tiempo}};

\draw[flechagruesa] (in) -- (f1);
\foreach \i [evaluate=\i as \j using int(\i+1)] in {1,...,5} \draw[flechagruesa] (f\i) -- (f\j);
\draw[flechagruesa] (f6) -- (db);

% rechazos: caída vertical bajo su filtro
\node[notaaviso, text width=2.6cm, align=center, minimum height=1.05cm] (des) at (2.35,-2.3)
  {\textbf{desconocido}\\{\tiny rechazo + registro}};
\node[notaaviso, text width=2.6cm, align=center, minimum height=1.05cm] (cua) at (7.05,-2.3)
  {\textbf{malformada}\\{\tiny cuarentena con motivo}};
\draw[flecha, color=ambar] (f1.south) -- (des.north);
\draw[flecha, color=ambar] (f3.south) -- (cua.north);
\end{tikzpicture}}
\caption{El pipeline de tratamiento de tramas, aporte central del trabajo. Cada filtro es
independiente y verificable por separado; las dos salidas de rechazo preservan la evidencia sin
contaminar las series.}
\label{fig:pipeline}
\end{figure}
